% !TEX program = xelatex

\documentclass[11pt, a4paper]{article}

% Required system preamble for fonts and language
\usepackage[a4paper, margin=2.5cm]{geometry}
\usepackage{fontspec}
\usepackage[english, bidi=basic, provide=*]{babel}
\babelprovide[import, onchar=ids fonts]{english}
\babelfont{rm}{Helvetica}

% Additional standard packages for formatting
\usepackage{xcolor}
\usepackage{tabularx}
\usepackage{booktabs}
\usepackage{enumitem}
\usepackage[colorlinks=true, linkcolor=blue!70!black, urlcolor=blue!70!black]{hyperref}

% Define some subtle colors for the document
\definecolor{primary}{HTML}{1A365D}
\definecolor{secondary}{HTML}{2B6CB0}

\begin{document}

% --- HEADER SECTION ---
\begin{center}
    {\huge \textbf{\textcolor{primary}{DFIR Investigation Report}}}\\[0.4cm]
    {\Large \textbf{Case:} [Insert Lab Name / Case ID]}\\[0.4cm]
    \begin{tabular}{r l | r l}
        \textbf{Analyst:} & [Your Name/Handle] & \textbf{Date:} & \today \\
        \textbf{Platform:} & CyberDefenders & \textbf{Status:} & Closed \\
    \end{tabular}
    \vspace{0.3cm}
    \hrule
\end{center}
\vspace{0.5cm}

% --- 1. EXECUTIVE SUMMARY ---
\section*{\textcolor{primary}{1. Executive Summary}}
The attacker performed targeted enumeration against the HTTP service hosted on the target IIS server. Subsequently, the threat actor initiated an SMB connection and successfully planted a web-accessible payload. A misconfiguration in the SMB shares allowed anonymous/unauthorized write access, enabling the attacker to drop an ASPX webshell (\texttt{shell.aspx}). This file contained the necessary code to establish a Command and Control (C2) communication channel. Upon execution by the server, the webshell facilitated the download of the final malicious payload.

% --- 2. METHODOLOGY & TOOLS ---
\section*{\textcolor{primary}{2. Methodology \& Tools Used}}
The following forensic tools were utilized to analyze network traffic, extract memory artifacts, and perform static analysis on the identified malware:
\begin{itemize}[label=\textcolor{secondary}{$\bullet$}, itemsep=2pt]
    \item \textbf{Wireshark:} Employed for deep packet inspection and network traffic analysis of the provided PCAP file.
    \item \textbf{Volatility 3:} Utilized for memory forensics to identify rogue processes and extract injected code.
    \item \textbf{PEStudio:} Used for static malware analysis of the extracted executables to identify strings, imports, and capabilities.
    \item \textbf{VirusTotal:} Leveraged for OSINT and signature verification of extracted file hashes.
\end{itemize}

% --- 3. INCIDENT TIMELINE ---
\section*{\textcolor{primary}{3. Incident Timeline}}
\textit{Note: All timestamps are recorded in UTC.}
\vspace{0.2cm}\\
\noindent
\begin{tabularx}{\textwidth}{@{} l l X @{}}
\toprule
\textbf{Date} & \textbf{Time} & \textbf{Event Description} \\
\midrule
2024-09-10 & 05:47:33 & First tree connect request observed to \texttt{10.0.2.15\textbackslash IPC\$} \\
2024-09-10 & 05:47:33 & Second tree connect request observed to \texttt{10.0.2.15\textbackslash Documents} \\
2024-09-10 & 05:48:52 & Payload (\texttt{shell.aspx}) Create Request initiated by attacker \\
2024-09-10 & 05:48:52 & Create Response generated by the server \\
2024-09-10 & 05:48:53 & Malicious payload (\texttt{shell.aspx}) successfully written to the server \\
2024-09-10 & 05:49:51 & Payload executed; outbound connection to C2 server established \\
\bottomrule
\end{tabularx}

\vspace{0.5cm}

% --- 4. TECHNICAL FINDINGS & ARTIFACTS ---
\section*{\textcolor{primary}{4. Technical Findings \& Artifacts}}

\textbf{Initial Access \& Execution:} \\
The investigation revealed that malware was pushed onto the target device utilizing an open SMB connection. The \texttt{Documents} share was misconfigured, lacking proper access controls, which allowed the threat actor to upload the \texttt{shell.aspx} file directly into a web-facing directory. 

\vspace{0.3cm}
\noindent \textbf{Indicators of Compromise (IoCs):}
\begin{itemize}[label=\textcolor{secondary}{$\bullet$}, itemsep=2pt]
    \item \textbf{Target IP Address:} 10.0.2.15
    \item \textbf{Attacker IP Address:} [Insert IP from PCAP]
    \item \textbf{C2 Domain / IP:} [Insert C2 destination from PCAP]
    \item \textbf{Dropped File:} \texttt{shell.aspx}
    \item \textbf{File Hash (SHA256):} [Insert hash you uploaded to VT]
\end{itemize}
\vspace{0.3cm}

% --- 5. MITRE ATT&CK MAPPING ---
\section*{\textcolor{primary}{5. MITRE ATT\&CK Mapping}}
The observed adversary behaviors have been mapped to the MITRE ATT\&CK framework to provide a standardized understanding of the attack lifecycle:

\vspace{0.2cm}\noindent
\begin{tabularx}{\textwidth}{@{} l l X @{}}
\toprule
\textbf{Tactic} & \textbf{Technique ID} & \textbf{Technique Name \& Context} \\
\midrule
Reconnaissance & T1595 & \textbf{Active Scanning:} Attacker enumerated the HTTP service on the target IIS server. \\
Initial Access & T1133 & \textbf{External Remote Services:} Exploited misconfigured, open SMB shares to gain a foothold. \\
Persistence & T1505.003 & \textbf{Web Shell:} Uploaded \texttt{shell.aspx} into a web-facing directory to execute commands. \\
Execution & T1059 & \textbf{Command and Scripting Interpreter:} Server executed the ASPX payload to initiate the next stage. \\
Command \& Control & T1071 & \textbf{Application Layer Protocol:} Payload established outbound C2 communication. \\
\bottomrule
\end{tabularx}
\vspace{0.5cm}

% --- 6. PERSONAL REFLECTION & SKILLS DEVELOPMENT ---
\section*{\textcolor{primary}{6. Personal Reflection \& Skills Development}}
\textit{As part of my continuous learning path in DFIR, this lab highlighted several key areas of growth that I am actively working to improve:}

\begin{itemize}[label=\textcolor{secondary}{$\bullet$}, itemsep=4pt]
    \item \textbf{Process Identification:} I recognized a need to improve my baseline knowledge of normal Windows OS processes. I am currently studying standard process trees to better identify malicious anomalies when using Volatility.
    \item \textbf{IIS \& ASPX Architecture:} To better understand webshell execution, I am dedicating time to learning how IIS processes ASPX files and handles backend execution.
    \item \textbf{Wireshark Efficiency:} Navigating the PCAP took longer than expected. I am currently building a custom Wireshark profile with specific color-coding rules for SMB and HTTP traffic to accelerate future triage.
\end{itemize}

\end{document}